\section{Secondary-data design}

The secondary-data design treats the project as a measurement reassessment of
Bard, Robertson, and Sorace \citep{BardEtAl1996}. It separates four kinds of
source from the observation channels introduced in
Section~\ref{sec:gramm-accept}: sources are routed to estimands, while channels
bear on grammatical inferences. Published summaries from Bard et al.\ can source
the historical claim, but can't support a direct replication without their
participant-level data.

Later judgment studies can support secondary analysis when they compare
elicitation methods on shared or comparable items and provide the rows needed
for the intended estimand. Computational benchmarks can show the afterlife of
acceptability labels, but they don't bear on magnitude estimation versus
Likert ratings. Construct-level discussion can separate acceptability from
grammaticality, naturalness, standardness, and correction.

\begin{table}[H]
\centering
\small
\begin{tabular}{@{}>{\raggedright\arraybackslash}p{0.19\linewidth}>{\raggedright\arraybackslash}p{0.22\linewidth}>{\raggedright\arraybackslash}p{0.25\linewidth}>{\raggedright\arraybackslash}p{0.24\linewidth}@{}}
\toprule
Source & Evidence level & Use in this paper & Unsupported claim \\
\midrule
Bard et al.\ 1996 & Published article & Historical target and scale-resolution claim & Direct replication without participant rows \\
Sprouse, Sch\"utze, and Almeida 2013; Sprouse and Almeida 2017 & Public row-level comparison data & Constrained secondary method comparison within approved reuse boundary & Raw-data redistribution or drift beyond the approved methodological scope \\
Langsford et al.\ 2018 & Published article and supplement & Reliability and response-style synthesis & New participant-level reliability model without raw rows \\
CoLA, BLiMP, SAD & Benchmark labels and documentation & Afterlife of formal acceptability measurement & Evidence about ME versus other human judgment methods \\
\bottomrule
\end{tabular}
\caption{Evidential routing before analysis. ME = magnitude estimation; SAD =
Syntactic Acceptability Dataset.}
\label{tab:evidence-routing}
\end{table}

The approved reuse boundary resolves the secondary-data side of the genre gate.
Public row-level data from the two comparison studies support a constrained
secondary analysis of method performance. The original studies had different
aims: \textcite{sprouse2013} estimated convergence between informal and formal
judgments, while \textcite{sprouse2017} studied task sensitivity and statistical
power. The present resolution analysis is a secondary use of their rows, not a
claim about what those studies were designed to test.

The verified public Langsford materials available to this project lack raw
participant responses. This limits Langsford et al.\ to the article-summary
level; the analysis doesn't attempt a participant-level
response-style or test-retest model across the literature.

The comparison files also require representation choices before modelling. For
the 2013 data, the primary forced-choice comparison uses the pair-level sign-test
file because forced choice is a pairwise task and that file records judgments
against expected patterns. The logistic file preserves item-level binary rows
and changes the unit of analysis, so it serves as a named sensitivity
specification.

For the 2017 data, yes/no rows with item IDs \texttt{B}, \texttt{G}, and
\texttt{M} are excluded from item-level comparisons.
They're absent from the materials workbook, occur in repeated early order
positions for every subject, and don't align with the magnitude-estimation,
Likert, or forced-choice item universe. This exclusion is a structural
inclusion rule, not an outcome-based trimming decision.

Each source is assigned to an estimand it can answer. The primary estimand is
the item-level acceptability signal
recoverable across methods. Secondary estimands are method-specific noise or
bias and decision agreement for theoretically relevant contrasts.

Participant response style and method-by-item instability are admissible only
when the dataset contains the participant-level rows and repeated structure
needed to estimate them. This ordering keeps the analysis from selecting
whichever scale looks best after the fact, the kind of forking path Gelman and
Loken warn about \citep{GelmanLoken2013,gelman2014statcrisis}.

The substantive comparison uses a measurement null, but a careful one. Magnitude
estimation, Likert ratings, forced choice, and Thurstonian choice are modelled
as readings of a common item-level acceptability signal through method-specific
response functions, not as the same signal plus method noise: the bounded
methods can compress the signal at the extremes, while the open ratio scale need
not. The single-dimension assumption is checked before any
method-specific term is read. With richer row-level evidence, that check would
compare a one-factor model against correlated-two-factor and bifactor
alternatives and inspect method residuals for item-level structure. In the
present reanalysis, the implemented counterpart is deliberately weaker:
a first-pass PCA/parallel-analysis check over aggregate item and contrast
matrices.

This ordering matters because the advantage Bard et al.\ actually claimed is
about resolution, recovering distinctions that fixed scales censor
\citep[35]{BardEtAl1996}. A plain common-signal-plus-noise null would book
that resolution as magnitude-estimation noise and let the method only lose; the
response functions and the dimensionality check let a resolution advantage
register instead. If one factor suffices, magnitude estimation still has to earn
its term; if it doesn't, the analysis turns to where the methods diverge.
Scrambled labels, method-label permutations, and simulations from the
response-function model would be pipeline controls rather than substantive
comparisons; the paper doesn't use them as evidence about ME.

After the chartered analyses were complete, a bounded multiverse tested whether
their conclusion depended on defensible representation choices. Its
specification universe was frozen before that analysis was run, but after the
primary outcomes were known, so it's a post-outcome robustness analysis rather
than a prespecified test. Three families vary ME and Likert scoring, endpoint
definitions and spread measures, raw- and rank-scale prediction, and independent
forced-choice or yes/no adjudication. The resulting 162 specifications test
scale-boundary spread, incremental out-of-sample information, and actual
contrast-sign disagreements without treating unexplained ME variance as an
advantage by itself. The complete decision rules and specification list are
versioned in \texttt{notes/sprouse-resolution-multiverse-plan.md}.
